\section{Progression remainders and extraction of primes}\label{sec:extraction}

\subsection{The complex progression remainder}

Retain fixed $M,b$ with $(b,M)=1$, $0<c_0<\eps/8$, and a smooth
nonnegative compact cutoff $\Phi$. We first prove the estimate needed
by the positive sieve.

\begin{proposition}\label{prop:complexremainder}
For the inert-array weights $H$ allowed in
\eqref{eq:unsigned-progression}, every $A>0$, and sufficiently large
fixed $N$,
\begin{equation}\label{eq:complexremainder}
 \sum_{\substack{r\leq x^{c_0}\\(r,M)=1}}
 \left|\sum_{\substack{r\mid u+1\\u+1\equiv b\ (M)}}
    \Phi(u/x)H(u)[f_z(u)T_x(u)-\Bx(u)]\right|
 \ll_A xL^{-A}.
\end{equation}
\end{proposition}

\begin{proof}
Here $H(u)=H(u_y)$ and is bounded for each fixed array collection.
Discard $u_y>D_0$ by the positive divisor envelope
\eqref{eq:rankin} and the summed rounding estimate
\eqref{eq:rounding}. In a genuine tuple term write
$u=\ell av$, where $a=u_y\leq D_0$ and $P^-(v)>y$.
The long product $\ell$ is a unit modulo $rM$ for large $x$.
Set $g=(b-1,M)$. Compatibility requires
\begin{equation}\label{eq:compatibility}
 (a,r)=1,\qquad (a,M)=g.
\end{equation}
When these conditions hold, $v$ lies in a reduced class modulo
$q=rM/g$. Prime by prime, partial removal of a prime-power factor
of $M$ fixes the valuation of $a$; full removal allows larger
valuations. In each case $a/g$ is a unit modulo $M/g$.
Roughness still excludes every prime at most $y$, including those
fully removed from the modulus. The model term $u=av$ has the same
compatibility and modulus, although its reduced residue can differ.

Insert \eqref{eq:bonferroni} in the genuine progression sum and in
its centered natural mean. Only divisors $e$ coprime to $q$ can
contribute in a reduced class. Their residual scales satisfy
$R/e\geq x^{\eps-o(1)}$, and the distinct moduli $rM/g$ are
less than $(R/e)^{1/3}$. Apply \eqref{eq:bvz}, then sum with the
$L^{O(1)}$ reciprocal divisor mass. The positive envelope controls
progression errors by CRT; natural-mean errors cost only a
logarithmic reciprocal-totient sum.

In the remaining natural means, the exclusion of primes $p\mid q$
with $p>y$ can be removed. After division by $\varphi(q)$, its
total cost is bounded using
\begin{align}\label{eq:largeq}
 \sum_{q\leq Mx^{c_0}}\frac1{\varphi(q)}
     \sum_{\substack{p\mid q\\p>y}}\frac1p
 &\leq\sum_{p>y}\frac1{p(p-1)}
           \sum_{t\leq Mx^{c_0}/p}\frac1{\varphi(t)}
 \ll L^{O(1)}/y.
\end{align}
We used $\varphi(pt)\geq(p-1)\varphi(t)$, including $p\mid t$.
This saves every logarithmic power. Lemma~\ref{lem:roughdensity}
at modulus one now gives, for fixed $a,\ell$, principal contribution
\begin{equation}\label{eq:presievemain}
 \frac{xz^df_z(a)H(a)}{a\ell\varphi(q)}
    \int\Phi(t)b_{z,y,N}(xt/(a\ell))\dd t.
\end{equation}

For the model, use unweighted rough counting in the reduced class
of $v$. Its density is
\begin{equation}\label{eq:roughq}
 \frac{V_y}{q\prod_{p\mid q, p\leq y}(1-1/p)}
 =\frac{V_y}{\varphi(q)}\prod_{p\mid q, p>y}(1-1/p).
\end{equation}
The $1/V_y$ in \eqref{eq:model} cancels $V_y$.
Equation~\eqref{eq:largeq} permits replacement of the final product
by one. The result is exactly \eqref{eq:presievemain}, including
every retained density coefficient. These terms cancel before
absolute values are taken. Sum the errors with the reciprocal
$a$-mass $L^{O(1)}$ and tuple mass $O(1)$, using
\eqref{eq:rounding} for the integer-counting errors.
\end{proof}

\subsection{A common positive sieve argument}

We now treat three first-endpoint differences:
\begin{equation}\label{eq:threecores}
 D_x(u)=
 \begin{cases}
 T_x(u)-C_x,&\text{unsigned progression law},\\
 \lambda(u)T_x(u),&\text{parity},\\
 f_z(u)T_x(u)-\Bx(u),&\text{nonreal mode}.
 \end{cases}
\end{equation}
Their summed progression remainders vanish to arbitrary logarithmic
precision by \eqref{eq:unsigned-progression},
\eqref{eq:signed-progression}, and
Proposition~\ref{prop:complexremainder}. For the unsigned case,
subtract $C_x$ times the remainder without the tuple.

Choose a large fixed even $h$, put $\gamma=c_0/(4h+3)$, and sieve
$u+1$ by primes at most $x^\gamma$ that do not divide $M$.
For the unsigned case apply the block sieve separately to
$\Phi\W T_x$ and $C_x\Phi\W$. For parity use the nonnegative
weights $\Phi\W T_x(1\pm\lambda)/2$.
For the complex case, $|\Bx|\leq1$ eventually, so with
$A_0=\Phi\W(T_x+1)$ the four weights
\[
 A_0\pm\Phi\W\Rea D_x,\qquad
 A_0\pm\Phi\W\Ima D_x
\]
are nonnegative. All these weights include the same progression
indicator. Each plus/minus pair has identical main terms and
arbitrary-precision summed remainders. The block sieve is applicable
because $x^{\gamma(4h+2)}<x^{c_0}$. Also
\[
 \frac1M\prod_{p\leq x^\gamma, p\nmid M}(1-1/p)
 =\frac1{\varphi(M)}\prod_{p\leq x^\gamma}(1-1/p)
\]
for large $x$. Subtracting the paired estimates gives
\begin{equation}\label{eq:sifted}
 \limsup_{x\to\infty}\frac{L}{xm_{\W}}
 \left|\sum_{\substack{u+1\equiv b\ (M)\\P^-(u+1)>x^\gamma}}
       \Phi(u/x)\W(u)D_x(u)\right|
 \ll_{M,d,\eps,\Phi}e^{-h}/\gamma.
\end{equation}
The surviving constant is independent of $K,h$ and the particular
fixed inert array. Array-dependent remainder constants vanish in
the first limit. In the complex case, the unsigned upper sieve and
\eqref{eq:modelsize} bound the presieved model by
$o(xm_{\W}/L)$ for fixed $h$ and array. Thus
\eqref{eq:sifted} also holds with $D_x$ replaced by $f_zT_x$.

\subsection{Replacing composite factors and retaining the coin coefficients}

A presieved composite in the support has a fixed bounded number
$j$ of prime factors, with $2\leq j\leq2/\gamma$ for large $x$.
Repeated large factors have negligible absolute mass.
Telescope each prime factor to the positive rough proxy of
Section~\ref{sec:framework}. The scalar criterion and
\eqref{eq:detphase} permit this for the unsigned, parity, and
complex statistics. Fixed characters impose the product residue.
Clipping, coefficient truncation, dyadic splitting, and logarithmic
separation have the source's original absolute error bounds.
Choose $K$ after these fixed losses and $\gamma$, independently of
the later inert exponents.

Write $v=\prod_{i=1}^jm_i$, $u=v-1$, and let $\mathcal C$ be the
full candidate set of small and big groups. A bar on $m_i$ removes
all candidate prime parts, including inactive groups. The stripped
coefficient is
\[
 \widetilde a(m_i)=\ind_{P^-(m_i)>W}
     \ind_{\bar m_i>x^\gamma}\kappa(\bar m_i),
 \qquad \kappa(t)=\frac1{V(W)\log t}.
\]
The source bounds its total positive tuple mass by
$O_{j,\gamma}(x/L)$ and the threshold/logarithm replacement error
by $o(x/L)$.

Conditional on the tuple, the source uses independent coins with
probabilities
\begin{equation}\label{eq:coins}
 s_g(u,v)=c_1q_0^{\omega_g(u)-1}\frac{\omega_g(u)}{V_g}
              (q_0/j)^{\omega_g(v)-1}\frac{\omega_g(v)}{V_g}.
\end{equation}
Choose $c_1>0$ so that these numbers lie in $[0,1]$; a zero hit
count gives zero. Let $\mathcal R$ be the success sets with at least
$c_2\log L$ groups of each kind. The source's success-probability
estimate says that the other outcomes have total positive proxy
mass $o(x/L)$. Discard them at this stage, before expanding signs.
All first-endpoint statistics under consideration are bounded.

The retained-event probability has the exact expansion
\begin{align}\label{eq:coinexpand}
 R(u,v)
 &=\sum_{S\in\mathcal R}\prod_{g\in S}s_g
                     \prod_{g\notin S}(1-s_g)\notag\\
 &=\sum_{S\in\mathcal R}\sum_{T\subseteq\mathcal C\setminus S}
                  (-1)^{|T|}\prod_{g\in S\cup T}s_g,
 \qquad 0\leq R(u,v)\leq1.
\end{align}
Each active set $P=S\cup T$ meets the hypotheses of the marked
correlation, and the number and coefficient mass of terms are
$L^{O(1)}$.

Outside candidate-prime-square exceptions, strip the active part
of $v$, obtaining $v_*$. For each factorization $e_1\cdots e_j=v_*$,
the active primes can be allocated independently to the $j$ slots.
The bars make the coefficients identical for all allocations.
The identity
\begin{equation}\label{eq:allocation}
 j^{\omega_P(v)}(q_0/j)^{\omega_P(v)-|P|}
       =j^{|P|}q_0^{\omega_P(v)-|P|}
\end{equation}
then gives the invariant core
\[
 G_P(v)=\sum_{e_1\cdots e_j=v_*}
 \prod_{i=1}^j\left[
  \ind_{P^-(e_i)>W}\ind_{\bar e_i>x^\gamma}\kappa(\bar e_i)\right].
\]
It satisfies $G_P(v)=G_P(v_*)$ and
$|G_P(v)|\leq L^{C_j}\tau(v_*)^{C_j}$.
Every expanded term has its actual coefficient
$(-1)^{|T|}c_1^{|P|}j^{|P|}$ multiplying the marked sum.
Candidate-square exceptions have exponentially small divisor-weighted
mass and remain negligible after this logarithmic expansion cost.

For the unsigned and parity cases, apply
\eqref{eq:unsigned-correlation} or
Lemma~\ref{lem:lambda-correlation}, with $r=b-1$ and
$\Psi(t)=t\Phi(t)$, and multiply by $x$.
Choose the accuracy after the expansion costs. This makes the
composite contribution $o(x/L)$.

For a nonreal mode, Proposition~\ref{prop:modeltransfer} replaces
the first core by $\Bx(u)$ with arbitrary logarithmic precision.
The principal contribution for a given $P$ is
\[
 \sum_{u+1\equiv b\ (M)}\Phi(u/x)\W(u)\Bx(u)
                W_{1,P}(u)W_{1,P}(u+1)G_P(u+1).
\]
Sum these terms with all their coefficients before estimating their
absolute values. Reversing \eqref{eq:allocation} and then
\eqref{eq:coinexpand} gives exactly
\begin{equation}\label{eq:aggregate}
 \sum_{m_1,\ldots,m_j}
 \ind_{v\equiv b\ (M)}\Phi((v-1)/x)\W(v-1)\Bx(v-1)
       R(v-1,v)\prod_i\widetilde a(m_i),
\end{equation}
up to the already controlled square exceptions.
The same $\Bx$ occurs for every active set, so it factors through
both identities. All other factors in \eqref{eq:aggregate} are
nonnegative, and $R\leq1$. Hence its absolute value is at most
\begin{equation}\label{eq:aggregatesmall}
 \sup_{u\asymp x}|\Bx(u)|\,O_{j,K,\gamma,\Phi}(x/L)=o(x/L).
\end{equation}
The source mass bound follows from $O(xV(W)^j)$ rough tuples per
product box, coefficient size $O_\gamma((LV(W))^{-j})$, and
$O(L^{j-1})$ boxes. The fixed inert weight costs a fixed constant.
Logarithmic expansion factors multiply only the arbitrary-precision
errors, not \eqref{eq:aggregatesmall}. This distinction permits roots
$z$ whose fixed positive exponent in \eqref{eq:modelsize} is small.

\subsection{Removing the inert marks}

We have bounded the prime sum with each fixed inert array by the
right side of \eqref{eq:sifted}. For any fixed integer $B\geq1$,
choose $B$ disjoint arrays of the same sufficient size $K$, and put
\[
 Z_B(u)=\left(\frac1B\sum_{a=1}^B\frac{\W_a(u)}{m_a}-1\right)^2.
\]
The source nonnegative deweighting estimate is
\begin{equation}\label{eq:deweight}
 \limsup_{x\to\infty}\frac Lx\sum_p\Phi((p-1)/x)Z_B(p-1)
 \ll C_K/B.
\end{equation}
It applies without a progression restriction. Restriction decreases
its left side. The other Cauchy factor is $O((x/L)^{1/2})$, because
$T_x-C_x$, $\lambda T_x$, and $f_zT_x$ are bounded on the support.
Consequently the unweighted statistics satisfy
\begin{equation}\label{eq:limsup}
 \limsup_{x\to\infty}\frac Lx
 \left|\sum_{p\equiv b\ (M)}\Phi((p-1)/x)F_x(p-1)\right|
 \leq C_{M,d,\eps,\Phi}\frac{e^{-h}}\gamma
                    +C_{d,\eps,\Phi}\sqrt{C_K/B},
\end{equation}
where $F_x$ is respectively $T_x-C_x$, $\lambda T_x$, or $f_zT_x$.
The model term has already been removed in the last case.

The order of limits is part of the proof. Fix the progression and
interior data, then $h,\gamma$, the finite range of $j$, and a
sufficient $K$. Fix any finite number $B$ of arrays. Choose the
physical families, high-time exponents, density order $N$, and error
precisions. Take $x\to\infty$ first, then $B\to\infty$, and finally
$h\to\infty$, choosing a new sufficient fixed $K$ for each $h$.
Since $\gamma^{-1}=(4h+3)/c_0$, \eqref{eq:limsup} tends to zero.
Constants depending on $N,K,B,\gamma$ in errors have disappeared in
the first limit.

\begin{theorem}[Interior estimates]\label{thm:interior}
Under Assumption~\ref{ass:inputs}, for every fixed reduced class
$b\pmod M$, every nonempty interior tuple statistic
\eqref{eq:tuple}, and every $\Phi\in C_c^\infty((0,\infty))$,
\begin{align}
 \sum_{p\equiv b\ (M)}\Phi((p-1)/x)(T_x(p-1)-C_x)&=o(x/L),
       \label{eq:interiorunsigned}\\
 \sum_{p\equiv b\ (M)}\Phi((p-1)/x)z^{\Omega(p-1)}T_x(p-1)&=o(x/L)
       \quad(z\ne1),\label{eq:interiorsigned}
\end{align}
where $z$ is any fixed root of unity different from one.
\end{theorem}

The proof used $\Phi\geq0$ for the positive sieve. A real smooth
compact cutoff is a difference of two such cutoffs, obtained with a
smooth dominating bump. Real and imaginary parts then give the
complex formulation by linearity.
